\ExerciseSection

\begin{enumerate}
\def\labelenumi{\arabic{enumi})}
\item
  Let \(f\) be a polynomial in \(n\) complex variables, with complex coefficients and not identically zero. Show that the complement in \(\mathbf{C}^n\) of the set \(S\) of points \(z = (z_i)\) such that \(f(z_1, z_2, \ldots, z_n) = 0\) (the ``algebraic variety'' with equation \(f = 0\) ) is connected. (If \(a, b\) are two distinct points of \(\mathbb{C}S\) , consider the intersections of \(S\) and the complex line through \(a\) and \(b\) .)
\item
  Let K be a non-discrete Hausdorff topological division ring (or field).
\end{enumerate}

\begin{enumerate}
\def\labelenumi{\alph{enumi})}
\tightlist
\item
  Let E be a left vector space of dimension n over K. If \((a_{i})_{1\leqslant i\leqslant n}\) is a basis of E, and if we transport to E the topology of \(K^{n}\) (the product of the topologies of the factors) by means of the bijective linear mapping
\end{enumerate}

\[
(x _ {i}) \rightarrow \sum_ {i = 1} ^ {n} x _ {i} a _ {i},
\]

then the topology so defined on E is independent of the basis \((a_{i})\) chosen, is compatible with the additive group structure of E, and is such that the mapping \((t, x) \to tx\) of \(K \times E\) into E is continuous. If F is a vector subspace of E, the topology induced on F by that of E is the same as the topology defined on F by starting from an arbitrary basis of F, as above; F is closed in E, and CF is dense in E unless F = E.

\begin{enumerate}
\def\labelenumi{\alph{enumi})}
\setcounter{enumi}{1}
\tightlist
\item
  Generalize Propositions 4, 5 and 6 of Chapter VI, § 1 to the division ring K. (To show that the mapping \(X \to X^{-1}\) is continuous in the neighbourhood of every nonsingular square matrix of order n over K, when K is not commutative, use induction on n: considering a neighbourhood of the unit matrix \(I_{n}\) of order n, note that every matrix X sufficiently near to \(I_{n}\) can be put in the form
\end{enumerate}

\[
\left( \begin{array}{c c c c} \mathrm{I} & \mathrm{O} & \ldots & \mathrm{O} \\ \lambda_ {2} & & & \\ \vdots & & I _ {n - 1} & \\ \lambda_ {n} & & & \end{array} \right) \left( \begin{array}{c c c c} \mu_ {1} & \mu_ {2} & \ldots & \mu_ {n} \\ \mathrm{o} & & & \\ \vdots & & Y & \\ \mathrm{o} & & & \end{array} \right),
\]

where \(I_{n-1}\) is the unit matrix of order n---i and Y is a nonsingular matrix of order n---i.)

\begin{enumerate}
\def\labelenumi{\alph{enumi})}
\setcounter{enumi}{2}
\item
  If K is commutative and E is an algebra of rank n over K, then the topology of E is compatible with its ring structure. Furthermore, if E has an identity element, the group G of units of E is open and dense in E, and the topology induced on G by that of E is compatible with the group structure of G. (If e is the identity element of E, and if x is a unit of E, then each of the equations xy = e, yx = e has one and one only solution y; for each of these equations consider the equivalent system of n linear equations which give the components of y with respect to a basis of E.)
\item
  If K is a non-complete field and E an algebra of rank n over K, and if the completion \(\hat{K}\) of K is a field, then the completion \(\hat{E}\) of the algebra E is that obtained by extending the ring of operators of E to \(\hat{K}\) . Hence construct examples of a topological division ring whose completion is not a division ring.
\end{enumerate}

\begin{enumerate}
\def\labelenumi{\arabic{enumi})}
\setcounter{enumi}{2}
\item
  Using Proposition 10 of Chapter I, § 3, no. 6, prove that the real projective space \(\mathbf{P}_{n}(\mathbf{R})\) is homeomorphic to the subspace of \(\mathbf{P}_{n}(\mathbf{C})\) consisting of these points which have at least one system of real homogeneous coordinates. \([\mathbf{R}_{n+1}^{*}\) being considered as embedded in \(\mathbf{C}_{n+1}^{*}\) , let A be a closed set in \(\mathbf{R}_{n+1}^{*}\) , saturated with respect to \(\Delta_{n}(\mathbf{R})\) ; let B be the set of all points \(\zeta x\) , where \(x \in A\) and \(\zeta \in U\) ; show that \(\overline{\mathbf{B}}\) is saturated with respect to \(\Delta_{n}(\mathbf{C})\) , and that A is the trace of \(\overline{\mathbf{B}}\) on \(\mathbf{R}_{n+1}^{*}\) .{]}
\item
  In \(\mathbf{P}_{n}(\mathbf{C})\) let \(H_{n}\) be the ``quadric'' defined by the equation
\end{enumerate}

\[
x _ {0} ^ {2} + x _ {1} ^ {2} + \dots + x _ {n} ^ {2} = 0.
\]

Show that every point of \(H_{n}\) has an open neighbourhood homeomorphic to \(C^{n-1}\) , that \(H_{n}\) is connected, and that the intersection of \(H_{n}\) with the complement of any complex projective hyperplane is connected (to prove this last assertion, reduce to the case \(n = 2\) ).

Show that \(H_{2}\) is homeomorphic to \(S_{2}\) , and \(H_{3}\) to \(S_{2} \times S_{2}\) (use the parametric representation of the quadric by means of its rectilinear generators).

\begin{enumerate}
\def\labelenumi{\arabic{enumi})}
\setcounter{enumi}{4}
\tightlist
\item
  The sphere \(\mathbf{S}_5\) being considered as a subspace of \(\mathbf{C}^3\) , consider the mapping of \(\mathbf{S}_5\) into \(\mathbf{R}^7\) which sends each point \(x = (x_1, x_2, x_3)\) of \(\mathbf{S}_5\) ( \(x_1, x_2, x_3\) complex) to the point \(y = (y_i)_{1 \leqslant i \leqslant 7}\) of \(\mathbf{R}^7\) , where
\end{enumerate}

\[
\begin{array}{c} y _ {1} = | x _ {1} | ^ {2} - | x _ {2} | ^ {2}, \quad y _ {2} = \mathcal {R} (x _ {1} \overline {{x}} _ {2}), \quad y _ {3} = \mathcal {I} (x _ {1} \overline {{x}} _ {2}), \quad y _ {4} = \mathcal {R} (x _ {1} \overline {{x}} _ {3}), \\ y _ {5} = \mathcal {I} (x _ {1} \overline {{x}} _ {3}), \quad y _ {6} = \mathcal {R} (x _ {2} \overline {{x}} _ {3}), \quad y _ {7} = \mathcal {I} (x _ {2} \overline {{x}} _ {3}). \end{array}
\]

This function has the same value at two points \(\pmb{x}, \pmb{x}'\) of \(\mathbf{S}_5\) such that \(\pmb{x}' = \zeta \pmb{x}\) , where \(|\zeta| = 1\) . Show that, by passing to the quotient, it induces a homeomorphism of \(\mathbf{P}_2(\mathbf{C})\) onto a subspace of \(\mathbf{R}^7\) .

¶ 6) Let K be a non-discrete Hausdorff topological division ring (or field).

\begin{enumerate}
\def\labelenumi{\alph{enumi})}
\tightlist
\item
  Endow the left projective space \(\mathbf{P}_n(\mathbf{K})\) with the quotient of the topology of \(\mathbf{K}_{n+1}^*\) by the equivalence relation \(\Delta_n(\mathbf{K})\) . Extend Propositions 1, 4 and 5 of Chapter VI, § 3, to \(\mathbf{P}_n(\mathbf{K})\) endowed with this topology. Show that \(\mathbf{P}_n(\mathbf{K})\) is connected if \(\mathbf{K}\) is connected, and otherwise is totally disconnected.
\end{enumerate}

* b) If K is locally compact (but not discrete), show that \(\mathbf{P}_{n}(\mathbf{K})\) is compact. (Let U be a compact neighbourhood of o in K; let \(a \in K\) be such that

\[
\lim _ {m \rightarrow \infty} a ^ {m} = 0;
\]

let S be the subset of \(K_{n+1}^{*}\) consisting of points \(\boldsymbol{x} = (x_{i})\) , all of whose coordinates \(x_{i}\) belong to U, and such that \(x_{k} \in a\) . \(\overline{U}\) for some index k; show that \(\mathbf{P}_{n}(\mathbf{K})\) is the canonical image of S). *

\begin{enumerate}
\def\labelenumi{\alph{enumi})}
\setcounter{enumi}{2}
\item
  Likewise, extend Definition 2 and Propositions 6 and 8 of Chapter VI, § 3 to the set \(\mathbf{P}_{n,p}(\mathbf{K})\) of projective linear varieties of \(p\) dimensions in \(\mathbf{P}_n(\mathbf{K})\) ; * also Proposition 7, assuming that \(\mathbf{K}\) is locally compact {[}for Proposition 7 consider, for each sequence \(\sigma\) of indices, the subset \(\mathrm{S}_{\sigma}\) of \(\mathrm{A}_{\sigma}\) consisting of matrices all of whose rows belong to the set S defined in \(b\) ); show that \(\mathbf{P}_{n,p}(\mathbf{K})\) is the canonical image of the union of the sets \(\mathrm{S}_{\sigma}]\) . * Generalize Proposition 9 of Chapter VI, § 3, when \(\mathbf{K}\) is a field.
\item
  If K is not commutative, let \(\mathbf{P}_{n,p}^{\prime}(\mathbf{K})\) denote the space of p-dimensional projective linear varieties in the right projective space of n dimensions over K {[}the topology of \(\mathbf{P}_{n,p}^{\prime}(\mathbf{K})\) being defined in the same way as that of \(\mathbf{P}_{n,p}(\mathbf{K})\) {]}. Show that \(\mathbf{P}_{n,p}(\mathbf{K})\) and \(\mathbf{P}_{n,n-p-1}^{\prime}(\mathbf{K})\) are homeomorphic {[}to every \((p+1)\) -dimensional vector subspace V of the left vector space \(E=K_{s}^{n+1}\) make correspond the orthogonal complement of V, which is a vector subspace of the dual \(E^{*}\) of E{]}.
\end{enumerate}

\begin{enumerate}
\def\labelenumi{\arabic{enumi})}
\setcounter{enumi}{6}
\item
  Extend Exercise 6 of Chapter VI, § 3 to spaces of projective linear varieties over C and H.
\item
  Extend Exercises 7, 8 and 9 of Chapter VI, § 3 to the spaces \(\mathbf{W}_{n,p}\) (no. 4).
\end{enumerate}

In the vector space \(H^{n}\) over the division ring of quaternions, we again denote by \(W_{n,p}\) the set of all sequences \((x_{k})_{1\leqslant k\leqslant p}\) of p vectors \(x_{k}=(x_{kj})_{1\leqslant j\leqslant n}\) such that

\[
\sum_ {j = 1} ^ {n} x _ {i j} \overline {{{x}}} _ {i j} = \mathrm{I} \quad (\mathrm{I} \leqslant i \leqslant p)
\]

and

\[
\sum_ {j = 1} ^ {n} x _ {i j} \bar {x} _ {i j} = 0 \quad (i \neq k)
\]

( \(\bar{x}\) denoting the conjugate of the quaternion \(x\) ). Extend Exercises 7, 8 and 9 of Chapter VI, § 3, to these spaces \(\mathbf{W}_{n,p}\) .

